\subsection{Propiedades de permanencia por paso al límite inductivo.}

En todo este número, $(A_\alpha,\varphi_{\beta\alpha})$ designa un \emph{sistema
inductivo de anillos} cuyo conjunto de índices es preordenado \emph{filtrante}; se pone
$A=\varinjlim A_\alpha$ y se designa por $\varphi_\alpha$ el homomorfismo
canónico $A_\alpha\to A$.

\begin{proposition}[5.13.1]
\label{IV.2.s5.13.1-fr}
\begin{enumerate}
\item[\textup{(i)}] Para todo índice $\alpha$, sea $\mathfrak a_\alpha$ un ideal de
$A_\alpha$ tal que, para $\alpha\leq\beta$, se tenga
$\varphi_{\beta\alpha}(\mathfrak a_\alpha)\subset\mathfrak a_\beta$.
Entonces los $\mathfrak a_\alpha$ forman un sistema inductivo para las
restricciones de los $\varphi_{\beta\alpha}$,
$\mathfrak a=\varinjlim\mathfrak a_\alpha$ se identifica canónicamente con un
ideal de $A$, y $A/\mathfrak a$ con
$\varinjlim(A_\alpha/\mathfrak a_\alpha)$. Si, además, se tiene
$\mathfrak a_\alpha=\varphi_{\beta\alpha}^{-1}(\mathfrak a_\beta)$ para
$\alpha\leq\beta$, entonces
$\mathfrak a_\alpha=\varphi_\alpha^{-1}(\mathfrak a)$ para todo $\alpha$.
\item[\textup{(ii)}] Recíprocamente, para todo ideal $\mathfrak a$ de $A$, si se pone
$\mathfrak a_\alpha=\varphi_\alpha^{-1}(\mathfrak a)$ para todo $\alpha$, los
$\mathfrak a_\alpha$ forman un sistema inductivo tal que
$\mathfrak a=\varinjlim\mathfrak a_\alpha$.
\end{enumerate}
\end{proposition}

Es claro que (ii) es un caso particular de (i). En (i), el hecho de que
$\mathfrak a$ se identifique canónicamente con un ideal de $A$, y
$A/\mathfrak a$ con $\varinjlim(A_\alpha/\mathfrak a_\alpha)$, proviene de la
exactitud del funtor $\varinjlim$ en la categoría de los grupos abelianos.
Además, $\mathfrak a$ es la unión de la familia creciente de los
$\varphi_\alpha(\mathfrak a_\alpha)$; por tanto, si
$x_\alpha\in A_\alpha$ es tal que
$\varphi_\alpha(x_\alpha)\in\mathfrak a$, existe $\beta\geq\alpha$ tal que
$\varphi_\alpha(x_\alpha)=\varphi_\beta(y_\beta)$ para un
$y_\beta\in\mathfrak a_\beta$, por tanto
$\varphi_\beta(\varphi_{\beta\alpha}(x_\alpha))=\varphi_\beta(y_\beta)$;
pero esto implica la existencia de un $\gamma\geq\beta$ tal que
$\varphi_{\gamma\beta}(\varphi_{\beta\alpha}(x_\alpha))=\varphi_{\gamma\beta}(y_\beta)$,
y, en consecuencia,
$\varphi_{\gamma\alpha}(x_\alpha)\in\mathfrak a_\gamma$, de donde
$x_\alpha\in\mathfrak a_\alpha$ si se tiene
$\mathfrak a_\alpha=\varphi_{\gamma\alpha}^{-1}(\mathfrak a_\gamma)$.

\begin{corollary}[5.13.2]
\label{IV.2.s5.13.2-fr}
Si $\mathfrak N_\alpha$ es el nilradical de $A_\alpha$, el nilradical de $A$
se identifica con $\varinjlim\mathfrak N_\alpha$, y $A_{\mathrm{red}}$ con
$\varinjlim(A_\alpha)_{\mathrm{red}}$. En particular, si los $A_\alpha$ son
reducidos, lo mismo ocurre con $A$.
\end{corollary}

\oldpage[IV-2]{132}

En efecto, es claro que
$\varphi_{\beta\alpha}(\mathfrak N_\alpha)\subset\mathfrak N_\beta$ para
$\alpha\leq\beta$, y que $\varinjlim\mathfrak N_\alpha$ es un ideal
nilpotente de $A$, por tanto contenido en el nilradical $\mathfrak N$ de $A$.
Recíprocamente, si $x\in\mathfrak N$, se tiene $x^h=0$ para un entero $h$; se
puede escribir $x=\varphi_\alpha(x_\alpha)$ para un índice $\alpha$ y un
$x_\alpha\in A_\alpha$; la relación
$\varphi_\alpha(x_\alpha^h)=0$ implica la existencia de un
$\beta\geq\alpha$ tal que
$\varphi_{\beta\alpha}(x_\alpha^h)=0$, o también $x_\beta^h=0$ poniendo
$x_\beta=\varphi_{\beta\alpha}(x_\alpha)$; por consiguiente,
$x_\beta\in\mathfrak N_\beta$ y, como $x=\varphi_\beta(x_\beta)$, esto
termina de probar el corolario, teniendo en cuenta (5.13.1).
